\paragraph{D.4f (Proved under the stated standard framework): the all-motives effectivity criterion detects the Tate conjecture.}
The remaining classification input in D.4b can be characterized
precisely. Fix $q$ and suppose that the category $\mathcal N_q$ of
pure numerical motives over $\mathbb F_q$ has the usual algebraic
Kunneth decomposition, is semisimple with finite-dimensional
rational morphism spaces, and has a faithful $\ell$-adic
realization factoring through numerical equivalence, for some
$\ell\ne p$. The last condition is essential: without it the
phrase \emph{eigenvalues of a numerical motive} has not been
justified by this realization.
Use the usual effective subcategory of direct summands of finite
sums of motives of smooth projective varieties. In this framework
the following are equivalent:
\begin{enumerate}
\item every motive whose Frobenius eigenvalues are algebraic
integers is effective;
\item every motive on which Frobenius is the identity is effective;
\item for every smooth projective $X/\mathbb F_q$ and every $r$,
the cycle map is onto the Frobenius-fixed subspace of
$H^{2r}(X_{\overline{\mathbb F}_q},\mathbb Q_\ell(r))$.
\end{enumerate}
Here (3) is the cycle-class form of the Tate conjecture on all
varieties over the fixed field; the framework already supplies
semisimplicity of the Frobenius actions used below.

\begin{proof}
The implication (1)$\Rightarrow$(2) is immediate. To prove
(2)$\Rightarrow$(3), first note that Frobenius is central in the
endomorphism algebra of every object: its realization commutes
with every correspondence, and the realization is faithful.
The centre of a finite-dimensional semisimple rational algebra
is a product of number fields. Frobenius consequently has a
square-free minimal polynomial. In particular, its fixed part
is cut out by a rational polynomial projector.

More explicitly, put $N=h^{2r}(X)(r)$. If its Frobenius minimal
polynomial has a factor $T-1$, write it as $(T-1)g(T)$ with
$g(1)\ne0$ and choose
\[
 P(T)\equiv1\pmod{T-1},\qquad P(T)\equiv0\pmod{g(T)}.
\]
The Chinese remainder theorem supplies $P\in\mathbb Q[T]$.
Let $M=\operatorname{im}P(F_N)$; if $T-1$ is absent, put $M=0$.
Then $F_M=1$, and its realization is exactly the fixed subspace
in (3). By (2), $M$ is effective.

Every effective pure motive of weight zero is Artin. Indeed
it is a summand of a finite sum of $h(Y)$ with no Tate
twists. The weight decomposition removes every $h^j(Y)$
with $j>0$, leaving only $h^0(Y)$; this last motive is that
of the finite etale scheme of geometric connected components.
For an Artin motive, a Frobenius element topologically
generates the Galois group of $\mathbb F_q$. Hence $F_M=1$
implies $M\simeq\mathbf1^{\oplus d}$.
The inclusions of these unit summands into $h^{2r}(X)(r)$
are rational algebraic cycles modulo numerical equivalence.
Their realizations form a basis of the fixed subspace.
This proves (3), including surjectivity after tensoring the
rational cycle space with $\mathbb Q_\ell$.

Conversely assume (3). Applied to products, it identifies
correspondences, followed by the given Kunneth projectors,
with all Frobenius-equivariant maps between realizations:
\[
 \operatorname{Hom}_{\mathcal N_q}(M,N)\otimes\mathbb Q_\ell
  \simeq
 \operatorname{Hom}_{F}(R_\ell M,R_\ell N).
\]
Injectivity after scalar extension follows from the
nondegenerate numerical trace pairing. Namely, choose
rational dual bases of the two correspondence spaces
$\operatorname{Hom}(M,N)$ and $\operatorname{Hom}(N,M)$
for composition followed by trace. A purported
$\mathbb Q_\ell$-linear relation between the realized
first basis, paired with the realized second basis,
has every coefficient zero. The realization preserves
these rational traces.
Surjectivity is exactly the Tate cycle statement on the
product in the appropriate degree and twist. Thus the
full-faithfulness input in D.4b holds. For an integral
eigenvalue of weight $i>0$, Honda--Tate applied to an
$i$th root supplies an effective abelian motive with that
eigenvalue. A common Frobenius constituent gives a nonzero
map from each simple constituent of the proposed motive
to an effective motive. The displayed isomorphism gives
such a map over $\mathbb Q$ as well: a nonzero tensor
product of a rational morphism space cannot arise from
the zero space. Semisimplicity splits it. Weight zero
and the exclusion of negative weight are as in D.4b.
This proves (1).
\end{proof}

The argument sharpens the remaining issue: even the
Frobenius-identity weight-zero case of the effectivity
claim would supply every Tate class in this framework.
The known abelian-type theorem D.4e cannot be extended
to all numerical motives merely by noting that their
eigenvalues are realizable on abelian varieties.

\paragraph{D.4g (Proved under the cited 2026 theorem): a common geometric polarization for arbitrary abelian families.}
The geometric preservation case can be enlarged beyond CM
families. Let $S$ be integral and of finite type over
$\mathbb Z$, and let $A_1,\ldots,A_t$ be abelian schemes
over $S$ equipped with relatively ample polarizations.
Let $\mathcal C$ be a pure Tate category generated under
direct sums, tensor operations and direct summands by
their degree-one motives and the Tate object, using
relative algebraic correspondences and projectors for
which exact faithful specialization tensor functors
\[
 F_s:\mathcal C\longrightarrow
          \operatorname{Mot}(\overline{k(s)};\mathrm{Ab})
\]
are defined and preserve the weight grading, the Tate
object and its fixed Tate structure. Thus cycles and projectors are assumed to
extend in the given family; a generic correspondence
with no specified specialization is not being used.
Assume Theorem I of \citet{D_VanHoftenSempliner2026}.
There is a common polarization $\Pi_{\mathcal C}$ such
that every $F_s$ preserves it. It is the unique
polarization containing the specified geometric forms
on the generating degree-one motives and the Tate
object. On the even subcategory, the distribution
assertion D.4a therefore holds with these geometric
polarizations and the fibre/polarization identification
specified there, whenever the chosen slope parities
exist.

\begin{proof}
By \citet[Theorem 3.3 and Corollary 3.4,
pp.~607--608]{D_Milne2002}, the CM Hodge hypothesis
supplied by the cited 2026 theorem gives the unique
geometric polarization $\Pi_s$ on the full category
of abelian motives over each algebraically closed
residue field. It also identifies numerical and
$\ell$-adic homological equivalence there. Theorem 2.1
ensures positivity of every geometric degree-one form.
The relatively ample forms in our statement specialize
to exactly these forms.

We give the tensor argument that ensures independence
of $s$. More generally, suppose an exact faithful tensor
functor $F:\mathcal C\to(\mathcal D,\Pi_{\mathcal D})$
carries fixed forms on a set of tensor generators to
positive forms. Work in each pure weight.
Tensor products, duals and direct sums
give a form on any expression in the generators.
If $X$ is a direct summand of that expression, restrict
the form to $X$. Its image under $F$ is a positive,
nondegenerate form, so the restricted form in
$\mathcal C$ is nondegenerate: faithfulness and
exactness reflect the invertibility of its associated
map to the dual. Its trace form on
$\operatorname{End}_{\mathcal C}(X)$ is positive
because this endomorphism space embeds in
$\operatorname{End}_{\mathcal D}(F X)$.

Thus every object has a Weil form whose image is
positive. Define the positive class on $X$ to be the
forms whose image belongs to $\Pi_{\mathcal D}(F X)$.
These classes are nonempty and satisfy the sum and
tensor compatibility axioms because their images do.
They form a polarization on $\mathcal C$; equivalently
this is the pullback construction implicit in
\citet[Lemma 1.2 and its preceding definition,
pp.~602--603]{D_Milne2002}.
Its classes are determined by the forms on the
generators, since the construction by sums, tensors,
duals and summands reaches every object.

Apply this argument to every $F_s$. In each case the
same relative generator forms and the same restricted
forms constructed from them are positive. Two
polarizations containing a common Weil form on every
object have the same compatibility classes. Hence the
pullbacks of all $\Pi_s$ coincide; call their common
value $\Pi_{\mathcal C}$.

The needed faithfulness of a specialization on this
specified category can also be checked geometrically.
A relative correspondence acts on the smooth proper
$\ell$-adic local systems. Vanishing at one geometric
point of connected $S$ means vanishing everywhere.
Numerical equivalence is homological equivalence on
the abelian fibres by Theorem 3.3, so quotienting by
that relation makes each specialization faithful.
For any two points of $S$, choose $\ell$ different
from their residue characteristics. Both lie in
the integral, hence connected, open $S[1/\ell]$,
where the preceding argument applies; the numerical
relation is independent of that choice.
The existence and typing
of the specified specialization functors remain part
of the statement, as they must for motives over a
base.

Finally restrict to even weights and apply D.4a.
After choosing an isomorphism to a representative
of the common real fibre class $\xi$, the same
polarization $\Pi_{\mathcal C}$ is transported at
every point. Changes of the chosen isomorphism give
the prescribed real inner-conjugacy ambiguity.
\end{proof}

This is a result for abelian-type families and their
extending algebraic correspondences. It does not
deduce from constancy of an unstructured torsor
class that an arbitrary choice of $\mu_2$-map is
geometrically positive. The requested assertion for
general geometric motives still needs precisely
that positivity and specialization compatibility,
in addition to a defined slope-parity map.
